\chapter{Единственный быстрый выход}
% полная версия: chapters/13_muscles.tex

Всё, что вы делаете в мире быстро, --- слово, взгляд, шаг --- это сокращение мышц. У машины есть и второй, медленный выход, химический, о нём глава шестнадцатая. А здесь --- мышцы.

Последнее звено цепочки --- \nt{АЦЕТ}-нейроны. Каждое мышечное волокно слушает ровно один такой нейрон, а один нейрон управляет группой волокон, от сотен до пары тысяч. Интересно, что соединение нерва с мышцей --- самое надёжное в теле. Внутри мозга соединения теряют больше половины сигналов, а здесь каждый щелчок выбрасывает в несколько раз больше вещества, чем нужно. На выходе ошибка стоит движения, и надёжность куплена с запасом.

\section*{Барабанная дробь становится толчком}

Один щелчок --- одно короткое подёргивание мышцы. Если щелчки идут редко, мышца дёргается отдельными рывками. Если часто, подёргивания сливаются в ровное усилие, как барабанная дробь сливается в гул. Мышца превращает поток импульсов в силу, как простейший цифро-аналоговый преобразователь.

\section*{Сила с плавающей точкой}

У мозга две ручки силы: как часто щёлкают нейроны и сколько групп волокон включено. Группы включаются по порядку --- от самых слабых к самым сильным, и никто этот порядок не вычисляет: маленькие нейроны просто легче возбуждаются. В итоге каждая новая группа добавляет примерно одну и ту же \emph{долю} к уже имеющейся силе. Когда вы берёте яйцо, силу дозируют крохотными порциями, когда поднимаете чемодан --- большими. Инженер узнает здесь число с плавающей точкой: порядок задаётся числом включённых групп, а точное значение --- частотой щелчков. Обратная сторона: чем сильнее движение, тем оно менее точно. По одной влиятельной гипотезе, поэтому наши движения такие плавные: плавная команда даёт наименьший разброс.

\section*{Тянуть, но не толкать}

Мышца умеет только тянуть. Поэтому каждый сустав обслуживает пара мышц, тянущих в разные стороны, --- снова два провода вместо знака. Разница их сил поворачивает сустав, а сумма делает его жёстче. Когда вы несёте полную чашку, вы напрягаете обе мышцы сразу: поворот тот же, но рука твёрже держит удар.

\section*{Обратная связь с опозданием}

В мышцах есть датчики растяжения, и самая короткая петля замыкается прямо в спинном мозге: мышцу растянули --- она сама сократилась, а мышца-противник на это время выключилась. Выключает её тормозной нейрон типа \nt{ГЛИЦ}: вот и нашлась работа этому типу.

Но у каждой петли есть задержка --- в спинном мозге десятки миллисекунд, через глаз --- сотни. А запоздалая обратная связь раскачивает систему, как шофёр, который рулит, глядя на дорогу с опозданием. Чем длиннее петля, тем медленнее она обязана исправлять. Граница, за которой спинномозговая петля начинает раскачиваться, --- около восьми колебаний в секунду. Это близко к частоте мелкой дрожи, которая есть у каждого здорового человека, и петля растяжения --- один из её вероятных источников.

Как же мы двигаемся быстро и точно? По распространённой гипотезе --- предсказываем: держим внутреннюю модель тела и рассчитываем результат команды, не дожидаясь датчиков.

\section*{Ритм без метронома}

Ходьба и дыхание ритмичны, но метроном для них не нужен. Две группы нейронов, тормозящие друг друга и постепенно устающие, работают по очереди сами: победитель устаёт, отдохнувший соперник берёт верх, и так снова и снова. Постоянная команда сверху --- «иди» --- превращается на месте в ритм «левой-правой».

\fullversion{13}
